\begin{trapbox}
\begin{description}[style=nextline, leftmargin=2.8em, labelsep=0.5em, itemsep=0.35em]
    \item[\textbf{\textcolor{CTrap}{易错点一（等号条件遗漏）}}] 只记住 $a=b=c$ 时等号成立，忽略两个相等且第三个为零的情况。
    \item[\textbf{\textcolor{CTrap}{正确理解（完整取等）：}}] 等号成立当且仅当 $a=b=c$ 或其中两个相等且第三个为 $0$。
    \item[\textbf{\textcolor{CTrap}{错因分析（记忆不全）：}}] 舒尔不等式的证明中，当两个变量相等且第三个为零时，表达式也为零，但容易被忽略。
\end{description}

\medskip
\begin{description}[style=nextline, leftmargin=2.8em, labelsep=0.5em, itemsep=0.35em]
    \item[\textbf{\textcolor{CTrap}{易错点二（范围忽视）}}] 在 $a, b, c$ 可能为负数或 $r \le 0$ 时直接使用舒尔不等式。
    \item[\textbf{\textcolor{CTrap}{正确理解（条件严格）：}}] 舒尔不等式要求 $a, b, c \ge 0$，$r > 0$；否则不等式不一定成立。
    \item[\textbf{\textcolor{CTrap}{错因分析（条件遗忘）：}}] 忽略定理的前提条件，导致在无效范围内应用结论。
\end{description}

\medskip
\begin{description}[style=nextline, leftmargin=2.8em, labelsep=0.5em, itemsep=0.35em]
    \item[\textbf{\textcolor{CTrap}{易错点三（形式记错）}}] 将舒尔不等式的符号记反，或记错项的次数与系数。
    \item[\textbf{\textcolor{CTrap}{正确理解（标准形式）：}}] 标准形式为 $a^{r}(a-b)(a-c) + b^{r}(b-a)(b-c) + c^{r}(c-a)(c-b) \ge 0$，当 $r=1$ 时可化为 $a^3+b^3+c^3+3abc \ge \sum_{\text{cyc}} a^2b$。
    \item[\textbf{\textcolor{CTrap}{错因分析（记忆混淆）：}}] 与其它不等式（如缪尔海德不等式）混淆，或只记了变形而忘了原始形式。
\end{description}
\end{trapbox}
